\chapter{总结与展望}
\section{总结}
本论文围绕基于ROS的多机器人协同搜索系统的设计与实现展开研究，针对传统单机器人系统在覆盖范围、响应速度和环境适应性方面的不足，设计了一种基于ROS2的分布式多机器人协同搜索系统。通过整合YOLOv5目标检测模型、传感器数据融合方案和改进的路径规划算法，构建了具备协作能力的多机器人搜索系统。主要涉及内容包括：
\begin{itemize}

    \item 多机器人协同架构：基于ROS2 Galactic的机器人操作系统，通过模块化设计支持多机器人实时数据交互。采用地图融合模块整合多机器人局部地图，构建全局环境模型，显著提升了系统的扩展性和鲁棒性。

    \item 目标检测优化：通过Gazebo仿真环境采集数据并训练轻量化YOLOv5s模型，结合时间戳同步策略和DBSCAN聚类算法，有效改善传感器数据延迟与坐标离散问题。实验表明，目标定位平均误差较小，且能过滤大部分误检干扰。

    \item 协同路径规划：提出基于A*算法与边界聚类的动态路径规划方法，结合障碍物膨胀策略和机器人避碰机制，实现了多机器人高效探索与协同避障。仿真测试中，双机器人系统在7分54秒内完成90平方米复杂环境的搜索任务，验证了算法的有效性。

    \item 全流程验证：通过Gazebo与RViz搭建仿真测试平台，完整验证了从环境建模、目标检测到协同搜索的全流程功能。系统在动态环境适应性和任务完成效率方面效果显著，为实际应用提供了可靠的技术参考。

\end{itemize}
\section{展望}

尽管本系统在仿真环境下取得了预期成果，但仍存在很多改进方向：
\begin{itemize}


    \item 实际场景适配：当前实验基于Gazebo仿真环境，物理引擎参数与真实环境存在差异。未来可在实体机器人平台上进一步验证，优化传感器噪声模型和运动控制精度。

    \item 算法性能提升：目标检测模块的误报率仍需降低，可通过引入多模态传感器（如热成像、深度相机）增强环境感知能力。路径规划算法可融合强化学习技术，实现动态任务分配与实时避障优化。

    \item 系统扩展性增强：现有架构支持双机器人协作，但大规模机器人集群的通信负载和决策效率尚未测试。后续可研究分布式共识算法等支持更多机器人协作的方案。

    \item 应用场景拓展：本系统目前针对较为简单的室内环境，未来可引入更多模块，扩展至安防巡检、灾害救援等复杂场景。同时，结合通信技术与物联网技术等，可进一步提升系统的实时性与部署灵活性。
\end{itemize}
本研究为多机器人协同搜索提供了一种经过仿真验证的技术方案，其模块化设计便于后续迭代。随着人工智能与机器人技术的深度融合，相信多机器人系统将在更多领域展现其智能化、集群化的优势，更好地被用于各种场景和需求。